% !TEX root = AImath.v5.tex
\chapter{人工智能的三种思维：从理论到应用的跨越}

人工智能的发展不仅代表着技术的进步，更标志着人类思维方式的演进。从最初的理论构想到如今的广泛应用，AI的每一次突破，都是不同思维模式交融与碰撞的结果。本章将带你梳理支撑AI发展的三种核心思维：数学思维、算法思维与工程思维，理解它们如何从理论、方法到实现，共同塑造了人工智能的基础。

\begin{introduction}[\faStar\;本章提要\;\faStar]
  \item 探讨人工智能发展的三种核心思维：数学思维、算法思维、工程思维；
  \item 分析三种思维在理论、方法与实现中的互补关系；
  \item 理解AI从抽象推理到现实应用的逻辑演进；
  \item  理解AI发展背后的三种思维模式；
  \item 能以数学、算法与工程的视角分析AI问题；
\end{introduction}

% 修改后版本：
\begin{introduction}[\faStar\;你将收获\;\faStar]
  \item 理解AI发展背后的三种思维模式；
  \item 能以数学、算法与工程的视角分析AI问题；
  \item 学会在理论与实践之间建立清晰的思维桥梁；
  \item 形成面向智能系统的综合思考框架。
\end{introduction}

人工智能的历史，既是一部探索“什么是智能”的追问史，也是一段不断突破技术与思维边界的演进历程。从“思维能否被计算”的哲学设想到今天的广泛应用，AI经历了无数次的试探、突破与挫折。每一次质疑都推动着思维向前，每一次停顿都孕育新的起点，每一次突破，又伴随着“智能”概念的重新定义。

换句话说，AI的发展不仅推动了技术的进步，也让人类不断反思自身的思维方式。正是这种“向外突破”与“向内反思”的互动，成为智能演化持续前进的内在动力。

纵观这一历程，我们可以看到，人工智能的发展并非一条平滑的直线，而更像是一种螺旋式的上升。理论的提出带动实践的验证，实践的反馈又催生新的思维与应用。正是在这种循环往复中，AI的进步从线性积累转化为“问题—方法—反思”的持续深化。

阅读本章时，不妨带着这一线索去观察，体会每个阶段中“问题—方法—反思”的交替节奏。思考人类如何在探索中，让抽象的思维逐渐转化为可计算、可实现的智能。以这样的视角回望，我们会更清晰地看到：范式在更替，思维在升级。
接下来，我们先用一个专门的小节勾勒出这三种思维的轮廓，然后再沿着AI发展的时间线，看看它们如何在不同阶段轮流登场、彼此交织。

\section{三种思维的融合：数学、算法与工程}

回顾人工智能的发展，我们会发现一个耐人寻味的规律：  
AI 的每一次飞跃，都源于思维方式的融合与重构。  
从符号主义的逻辑推理，到连接主义的神经计算，再到深度学习的经验学习——这些阶段并非彼此取代，而是思维范式的层层叠加，共同构筑出今日AI的广阔图景。

在本节中，我们先暂时放下具体的历史事件，用“三种思维”的视角搭起一个观察框架：数学思维、算法思维与工程思维。后面再回到AI的发展历程中，你可以一边看故事，一边练习用这三种“镜头”去解读。

\subsection{数学思维：形式化与可计算}

数学思维为AI提供了最坚实的基础——形式化的语言与推理工具。  
无论是符号逻辑、概率论，还是优化理论，数学让我们能够用严谨的方式去描述智能问题，使系统得以被定义、分析与验证。  
数学思维的真正力量在于：它能把模糊的直觉，转化为可以计算、可以推理的结构。

举个例子，在机器学习中，我们用概率模型来刻画不确定性，用优化方法去寻找最优解，用矩阵与向量表示高维数据。  
这些数学工具不仅构成了算法的骨架，更塑造了一种理解世界的方式。  
正是数学，使AI从“经验的堆积”转向“逻辑的生长”，从而拥有了内在的自洽性。

\subsection{算法思维：结构化与过程化}

如果说数学思维回答的是“能不能算”，  
那么算法思维则回答“怎么去算”。  
它强调结构与过程，关注如何在有限的资源与时间中高效实现目标。  
算法不仅是数学公式的执行者，更是一种组织思维、拆解问题的策略。

算法思维让AI具备了真正的“可操作性”。  
例如，在搜索问题中，我们通过启发式算法缩小搜索空间；  
在学习问题中，用梯度下降法优化模型参数；  
在推理问题中，借助动态规划或图算法实现高效求解。  
算法，让抽象的数学模型变成了可以运行的程序，也让智能的想法第一次有了实践的路径。

\subsection{工程思维：实现与扩展}

工程思维关注的是“如何让算法真正落地”。  
它不仅关乎软件实现，还涉及硬件优化、系统架构与数据管理等复杂问题。  
工程思维强调规模化、稳定性与鲁棒性，是AI从研究原型走向真实世界的关键桥梁。

比如，一个深度学习模型在实验室里也许只需几张 GPU 就能运行；  
但在工业级应用中，还要面对数据存储、分布式计算、网络通信与功耗优化等现实挑战。  
只有具备工程思维，我们才能让AI系统在复杂环境中稳定、持续地发挥作用。

\subsection{三种思维的交织与启示}

数学思维、算法思维与工程思维，并非彼此独立的三个世界。  
它们共同构成了AI的认知框架：数学提供原理，算法提供方法，工程提供实现。  
AI的发展史，正是一部三种思维不断交织、相互滋养的演化史。

这种融合，让AI完成了从理论到实践、从模拟到创造的跨越。  
从形式化的推理，到可计算的模型，再到可应用的系统，  
每一次进步，都是思维方式拓展的结果。  
而这种融合，不仅推动了AI本身的发展，也让我们得以用新的视角去理解“智能”这个古老而永恒的问题。

在接下来的章节中，我们就带着这三种思维，一起重新走一遍人工智能的发展历程：看看每一个关键阶段，究竟是哪几种思维在主导，哪几种思维在补位。

\section{AI的萌芽期：理论起源与计算基础}

人工智能的萌芽可以追溯到古代哲学家对思维本质的思考，但它真正进入科学范畴，是在20世纪上半叶。这一时期，数学、逻辑学与计算机科学的兴起，为AI奠定了坚实的理论根基；而计算机技术的突破，则让“思维的机械化”从设想变为可能。

如果用本章的视角来看，这一阶段主要是\emph{数学思维}和\emph{工程思维}在铺路：前者负责把“思维”形式化，后者负责把“计算”做强大。算法思维此时还只是隐约可见的轮廓。

\subsection{理论起源：从哲学思辨到科学构想}

“智能的本质是什么？”——这个问题困扰人类已久。从古希腊的哲学思辨到现代科学的理性研究，人类对智能的理解，经历了从感性直觉到形式化分析的漫长转变。

早在古希腊时期，亚里士多德便提出了“形式逻辑”的概念，试图用规则来描述推理的过程。他在《工具论》中系统阐述了三段论推理，这种把思维过程形式化的尝试，为后来的人工智能奠定了哲学基础。亚里士多德的贡献在于，他第一次尝试把人类的推理活动抽象为可操作的规则——这正是现代AI符号推理的思想雏形。

17世纪，莱布尼茨提出了“通用符号”（characteristica universalis）的概念，设想用符号与规则来表达所有知识与推理。他梦想创造一种“推理演算”（calculus ratiocinator），通过机械化的符号操作解决一切理性问题。这一宏大的设想，直接启发了后来的符号主义AI研究。可以说，莱布尼茨是人工智能最早的哲学先驱之一。

19世纪，数学逻辑的进步为AI奠定了更加坚实的理论基础。布尔代数为逻辑推理提供了精确的数学工具，而弗雷格的谓词逻辑进一步完善了形式化推理的框架。正是这些工具的成熟，使得人类第一次能够设想：思维的过程可以被计算机形式化地再现，从而为20世纪AI的诞生铺平了道路。

\subsection{计算机的诞生：从算盘到电子计算机}

如果说逻辑学为AI提供了思想的基石，那么计算工具的发展，则为AI提供了落地的物质支撑。从古代的算盘到现代的电子计算机，每一次计算工具的革新，都在推动AI理论的演进，也让“智能的机械化”逐渐变得可行。

计算工具的发展史，本身就是人类追求“智能机械化”的缩影。17世纪，帕斯卡发明了第一台机械计算器——帕斯卡计算器，能够完成加减运算。虽然功能简单，却传递出一个关键思想：部分智能活动（如计算）可以由机器承担。莱布尼茨在此基础上改进出可乘除的步进计算器，并提出了“二进制计算”的设想，为后来的计算机语言奠定基础。19世纪，查尔斯·巴贝奇设计了“分析机”（Analytical Engine），首次具备了存储器、运算器与控制器等现代计算机的核心结构。而他的合作者阿达·洛夫莱斯在编程过程中提出：机器或许能拥有创造力——这也许是最早的AI思想萌芽。

电子技术的突破，为AI的实现奠定了真正的技术基石。20世纪40年代，电子管技术的成熟让“电子计算机”成为现实。1946年，ENIAC——电子数值积分计算机——在宾夕法尼亚大学诞生，宣告电子计算机时代的到来。这台重达30吨的庞然大物虽然能耗惊人，但其运算速度远超以往的机械计算器，为复杂数值计算打开了大门。更重要的是，ENIAC具备可编程性：通过重新连接电缆与开关，它能执行不同任务。这种灵活的结构，为通用计算机的出现埋下了伏笔，也成为实现通用人工智能的关键特征。

冯·诺依曼架构的确立，为智能系统的发展奠定了核心理论。1945年，约翰·冯·诺依曼在《关于EDVAC的报告草案》中提出“存储程序计算机”概念——让程序与数据共存于同一存储器中，计算机因此能够修改自己的程序。这一设计不仅提升了计算机的灵活性，也为AI的发展提供了新的思维方式。其核心思想在于：程序本身也可以被计算和操控。这种“自指性”（self-reference）启发了后来的机器学习与自适应系统，使机器第一次具备了“自我修改”的潜力。冯·诺依曼本人在《计算机与人脑》中更进一步，比较了人脑与计算机的原理，提出了关于机器智能的诸多洞见，为AI研究奠定了理论坐标。

从本节看，数学思维让“智能”第一次有了可形式化的样子，工程思维让“计算”具备了足够的力量。接下来，算法思维将登场，把二者串联成真正可执行的“智能过程”。

\section{1950-1970年代：AI的奠基与早期探索}

20世纪50年代，被普遍认为是人工智能真正诞生的起点。此时，计算机技术日趋成熟，为AI研究提供了现实的平台；同时，一批杰出的科学家为这一新兴领域奠定了理论根基与研究方向。从图灵的理论构想到达特茅斯会议的召开，AI正从哲学思辨走向可验证的科学探索。

在这一时期，\emph{数学思维}给出了“可计算性”的理论边界，\emph{算法思维}开始以具体程序的形式出现，而\emph{工程思维}则通过早期计算机系统为这些想法提供实验场。

\subsection{图灵的贡献：从计算理论到智能测试}

阿兰·图灵被视为人工智能理论的奠基者之一。他不仅建立了计算理论的基础，更重要的是，他深入思考了“机器是否能够思考”这一核心问题。图灵的思想，为AI研究提供了理论框架与哲学方向，至今仍在影响这一领域的发展。

1936年，图灵提出了著名的“图灵机”概念——一个抽象的数学模型，用以精确定义“可计算”的边界。图灵机结构简单：只有一条无限长的纸带、一个读写头和状态控制器，却能模拟任何算法的执行过程。这一模型的深远意义在于，它证明了所有可计算问题都能以统一的形式描述与求解。这为后来的计算机科学奠定了理论框架，也为AI提供了理解“计算与智能”关系的基础。

更重要的是，图灵机的提出引出了一个极具启发的问题：如果人类的思维过程能够被算法化，那么机器是否也能拥有与人类相当的智能？这一思想为AI研究提供了理论依据，也开启了人类关于“机器思维可能性”的哲学讨论。

1950年，图灵发表了具有里程碑意义的论文《计算机器与智能》（Computing Machinery and Intelligence），正式提出了“图灵测试”这一著名思想。图灵设想：如果人与机器通过文字对话，而人类无法分辨对方究竟是人还是程序，那么这台机器就可以被视为“具有智能”。这一测试，为判断机器智能提供了最早且最具影响力的标准。

图灵测试的重要性，并不仅仅在于它提供了一种衡量机器智能的方式。更深层的意义在于：它体现了图灵对智能本质的洞见——智能不必复制人类的思维过程，而在于能否表现出与人类相当的行为。这种“行为主义”视角，为后来的AI研究指明了方向，也让研究者避免陷入对“意识”与“理解”的无休止争论。

在论文的结尾，图灵甚至大胆预言：到2000年，计算机将能够通过图灵测试。尽管这一目标尚未完全实现，但这份远见依然令人敬佩。图灵不仅看到了AI未来的可能，也预见了它在发展中必然面对的技术与伦理争议。

\subsection{达特茅斯会议：AI学科的正式诞生}

1956年夏，美国达特茅斯学院举办了一场具有历史意义的学术会议，被视为人工智能作为独立学科的正式诞生。这次为期两个月的会议由约翰·麦卡锡、马文·明斯基、克劳德·香农、艾伦·纽厄尔、赫伯特·西蒙与纳撒尼尔·罗切斯特等人共同组织。会议首次使用了“人工智能”这一术语，也确立了AI作为一门科学研究领域的身份。

达特茅斯会议的召开并非偶然，而是20世纪50年代科学技术发展到一定阶段的必然产物。二战后的技术浪潮中，计算机迅速崛起，控制论与信息论蓬勃发展，为AI研究提供了肥沃的土壤。与此同时，一批科学家开始认真思考“机器能否学习”的问题，迫切需要一个交流思想与整合资源的平台。会议的组织者在申请书中写下这样一句话：“我们假设，学习的每一个方面，或智能的任何特征，都能够被精确描述，使机器可以对其进行模拟。”这段文字充分展现了早期AI研究者的雄心与信念。

会议的参会者来自不同学科，奠定了AI研究的多元化基础。他们分别在语言、逻辑与信息处理等领域做出了重要贡献。  
约翰·麦卡锡不仅创造了“人工智能”一词，还发明了LISP语言，并提出“常识推理”和“情境演算”等概念。  
马文·明斯基在感知器理论与知识表示方面贡献突出，其《心智社会》提出：心智由多个简单智能体协作而成。  
克劳德·香农为AI提供了信息论的数学支撑，他设计的国际象棋程序成为后续博弈AI的启蒙。  
艾伦·纽厄尔与赫伯特·西蒙共同开发“逻辑理论机”，提出“物理符号系统假设”，奠定符号主义AI的理论框架。  
纳撒尼尔·罗切斯特则在神经网络与早期机器学习上迈出了先行步，为联结主义AI埋下种子。

达特茅斯会议的意义，远不止一次普通的学术聚会。  
首先，它确立了AI研究的基本目标——通过计算机程序模拟人类智能，为后续研究提供了清晰方向。  
其次，它展现了早期AI研究者的理想主义精神：他们相信机器不仅能理解世界，还能逐渐接近甚至超越人类的智能，这种信念成为AI早期发展的驱动力。  
再次，会议确立了AI研究的跨学科传统，来自数学、计算机科学、心理学与神经科学的思想在此汇流。  
最后，它提出了一个延续至今的方法论：用程序去验证理论，让理论在实践中生长。

\subsection{符号主义AI的早期发展}

达特茅斯会议之后，AI研究主要沿着“符号主义”的路线展开。符号主义认为，智能的核心在于符号操作——只要能定义清晰的符号及其规则，机器就能表现出智能行为。这一思想在AI发展的早期占据主导地位，并取得了一系列令人瞩目的成果。

“逻辑理论机”（Logic Theorist）是符号主义AI的第一个里程碑，也是AI历史上第一个真正意义上的智能程序。由纽厄尔、西蒙与肖（J.C. Shaw）共同开发，它能够证明《数学原理》（Principia Mathematica）中的38个定理，其中有一个证明甚至比原书更为简洁。这一成果极大地鼓舞了研究者的信心，展示了机器也能进行复杂的逻辑推理。其核心思想，是将问题转化为符号结构，通过符号操作完成求解。程序采用“启发式搜索”，在可能的证明路径中选择最有希望的方向——这一方法也成为此后AI问题求解的基本范式。

“通用问题求解器”（General Problem Solver，简称GPS）是另一个重要的早期AI系统。它试图通过“手段—目的分析”（means-ends analysis）来解决各种问题，体现了符号主义AI对“通用智能”的理想追求。其核心思想是：将问题表示为初始状态与目标状态，并通过一系列操作使二者逐步接近。  
在处理简单任务时，GPS的表现令人印象深刻，但当面对复杂的现实问题时，它的局限便显现出来——这些问题往往需要庞大的领域知识与更复杂的推理过程。正因如此，研究者开始重新思考“知识表示”和“领域专门化”，这也为后来的专家系统发展铺平了道路。

约翰·麦卡锡发明的 LISP（LISt Processing）语言，为符号主义AI提供了强大的编程工具。这门语言的设计哲学，深刻体现了符号主义AI的核心理念。  
LISP最独特的特征在于，它将程序与数据都表示为列表结构。这种“同构性”（homoiconicity）让程序能够操作自身，从而具备了自适应与学习的潜力。一个程序可以生成或修改另一个程序——这种灵活性，为AI研究带来了前所未有的表达能力。  
此外，LISP以符号处理为核心，与当时主要用于数值计算的 FORTRAN 等语言不同。它提供了丰富的列表操作函数，使复杂符号结构的构造与操作变得直观而高效。其函数式编程范式尤其适合处理递归定义的结构，如语法树与知识图谱。正因如此，LISP长期成为AI研究的首选语言，并在部分领域至今仍发挥着独特作用。

从“三种思维”的角度看，这一阶段是数学思维和算法思维的高光时刻：逻辑与符号提供了严整的结构，程序与搜索提供了可执行的过程。但工程思维的约束尚不明显，现实世界的复杂性也还没有真正登场。

\section{黄金期与第一次寒冬：理想与现实的碰撞}

1960年代常被称为人工智能的第一个“黄金时代”。这一时期，AI研究成果层出不穷，研究者的乐观情绪达到了前所未有的高度。然而，随着问题复杂度的提升与计算资源的局限显现，理想与现实之间的落差逐渐暴露。进入1970年代，AI迎来了第一次重大挫折——这场“寒冬”既带来停滞，也留下了深刻的启示，为后续发展奠定了反思的基础。

可以理解为：符号主义主导下的数学思维与算法思维，走到了一个“几乎无所不能”的乐观顶点，随即在现实世界的检验中遭遇工程与复杂性的反击。

\subsection{符号主义AI的黄金时代}

在这段黄金时期，符号主义AI迎来了最辉煌的阶段。无论是自然语言处理还是早期机器学习，都出现了突破性成果，似乎印证了早期研究者的乐观预期——智能机器的梦想，仿佛触手可及。

ELIZA 是这一时期最具代表性的自然语言处理程序之一，也是第一个引起公众广泛关注的 AI 系统。它由约瑟夫·魏岑鲍姆（Joseph Weizenbaum）在 MIT 开发，能够与用户进行简单的文字对话，模拟心理治疗师的语气与反应。  

从技术上看，ELIZA 的原理并不复杂：主要依靠“模式匹配”和“模板替换”。程序预设若干对话模板，当用户输入语句时，系统会寻找匹配的关键词并生成相应回复。例如，当用户说“我感到沮丧”时，ELIZA 可能回应：“你为什么感到沮丧？”——看似体贴，实则程式化。  

然而，正是这种简单的互动产生了惊人的心理效果。许多用户在与 ELIZA 交谈时表现出强烈的情感投入，甚至相信它真正“理解”了自己。这一现象后来被称为“ELIZA 效应”，提醒我们：人类极易把智能归因于具备表面理解能力的系统。

与 ELIZA 不同，Terry Winograd 在 MIT 开发的 SHRDLU 进一步迈向了“理解与行动”的结合。它能理解关于“积木世界”的自然语言指令，并在虚拟空间中执行操作。例如，当用户发出“把红色的积木放在蓝色的积木上”这样的指令时，SHRDLU 不仅能“理解”语义，还能完成动作。  
这一系统的意义在于：它把语言理解与问题求解联系起来，让人们看到了机器理解语言、执行任务的希望。但 SHRDLU 的局限也同样明显——它只能在“积木世界”这样受控的封闭环境中运行，无法面对真实世界的复杂与模糊。这也预示着符号主义AI面临的核心难题：如何从“实验室的智能”走向“世界的智能”。

尽管当时的AI研究以符号推理为主，但“机器能否自己学习”的想法开始萌芽。研究者逐渐意识到：手工编写所有规则几乎不可能，机器必须具备自我学习的能力。  
1957年，弗兰克·罗森布拉特（Frank Rosenblatt）提出了“感知器”（Perceptron）模型——一种简单的神经网络，可以通过训练学习线性分类任务。罗森布拉特还证明了感知器算法的收敛性：对于线性可分问题，它总能在有限步内找到正确的分类边界。这一成果让人们首次看到了“机器学习”的可行性。  
然而，1969年，马文·明斯基和西摩·帕普特在《感知器》一书中指出了它的致命局限：单层感知器无法处理非线性可分问题（如异或问题）。这场批评让神经网络研究陷入低潮，也标志着AI历史上第一次理想与现实的碰撞。

\subsection{自我学习的探索}

在早期研究中，AI学者逐渐意识到：仅靠手工编程去穷尽所有知识与规则并不现实。机器若要具备真正的智能，必须能够从经验中自我学习。这一认识推动了早期机器学习的萌芽——虽然技术尚处在原始阶段，却为后来的发展奠定了关键基础。

“概念学习”是早期机器学习的重要方向。研究者希望开发出能像人一样，从实例中归纳概念的算法。帕特里克·温斯顿（Patrick Winston）设计的“积木学习程序”正是其中的代表。它能够从正例与反例中学习出“拱门”等结构概念。  
程序通过分析正例的共同特征与反例的差异，逐步形成对概念的理解。例如，系统观察多个拱门后，会总结出：拱门由两个支柱和一个横梁组成，且支柱必须支撑横梁。  
这一研究揭示了学习的复杂性：即使是“拱门”这样看似简单的概念，也包含丰富的结构关系与逻辑约束。正因如此，概念学习算法通常只能在受限领域内有效，难以直接迁移到真实世界的多样与复杂。

与“演绎”相反，“归纳推理”从特殊到一般——它让系统能从具体经验中提炼出普遍规律。对于机器学习而言，这种推理方式具有关键意义：它使机器能从有限样本中获得可泛化的知识。  
早期的归纳学习算法多基于符号表示，试图从训练实例中归纳逻辑规则。例如，给定若干事实：“鸟会飞”、“企鹅是鸟”、“企鹅不会飞”，算法可能归纳出更一般的规则——“鸟会飞，除非它是企鹅”。  
然而，这项研究也揭示了学习的一个根本难题——归纳偏置（inductive bias）。从有限数据中可以推导出无数可能规律，因此学习系统必须带有某种偏向，以决定哪一种规律更“合理”。这种偏向可以来自先验知识、简单性原则，或某种启发式选择。

总体来看，这一时期的机器学习研究在理论上取得了重要突破，但距离广泛应用仍有距离。大多数算法只能在简化环境中运行，难以应对现实世界的复杂与不确定。这些局限，虽使AI一度陷入低谷，却也为后来的理论深化和算法创新提供了方向与动力。

\subsection{第一次“AI寒冬”的到来}

1970年代初，人工智能研究迎来了第一次重大挫折——后来被称为“AI寒冬”。  
这一时期的低谷并非偶然，它源于多方面的因素：一方面是技术能力的限制，另一方面则是人们对智能机器寄予了过高的期望。  
尽管“寒冬”带来了停滞与失落，但它也迫使研究者重新审视AI的本质与边界，从而孕育出新的思想萌芽。

随着问题规模的扩大，AI系统迅速暴露出计算复杂性的瓶颈。许多表面上看似简单的任务——如机器翻译与图像识别——实际上拥有极高的计算需求，而当时的硬件远不足以支撑。  
以机器翻译为例，早期研究者曾认为，只要建立词汇对照与语法规则，就能让机器实现自动翻译。然而，真正的翻译过程还涉及语义理解、上下文关联乃至文化背景，远比想象复杂。  
1966年，美国科学院发布的 ALPAC 报告指出，机器翻译的质量远不及人工翻译，且成本高昂、效率低下。此报告直接导致政府削减研究经费，使机器翻译陷入停滞。  
图像识别也面临相似困境。虽然人类轻松识别图像中的物体，但机器要做到同样的事情，却需经过特征提取、模式匹配和上下文理解等多个步骤，每个环节都带来巨大的计算负担。当时的算法与硬件都无法支撑这种复杂性。

符号主义AI的另一大瓶颈在于“知识表示”。  
它依赖于清晰、完备的逻辑结构，但现实世界的知识往往模糊、片面且充满矛盾。如何在计算机中表示这种复杂性，成为难以逾越的障碍。  
尤其是“常识知识”的表示问题更为突出。人类日常掌握着大量看似简单的常识——如“水往低处流”、“人需要呼吸”——但这些知识往往隐含在语境中，难以形式化地表达。  
更复杂的是，现实中的知识常伴随例外与不一致，例如“鸟会飞”这一规则有诸多例外（如企鹅、鸵鸟）。如何既保持知识体系的整体一致性，又能灵活应对例外，成为AI知识表示长期难解的课题。

早期的乐观预期未能兑现，AI研究的资金支持随之急剧减少。政府与企业的投资热情下降，许多研究项目被迫中止。  
1973年，英国政府委托詹姆斯·莱特希尔（James Lighthill）撰写评估报告。报告严厉批评AI研究未能实现承诺，建议削减经费——这一“莱特希尔报告”几乎使英国AI研究陷入停顿。  
美国的情况也相似。国防部高级研究计划局（DARPA）作为主要资助者，因AI项目未达预期而减少投资，迫使大批研究者转向其他领域。AI由此进入了第一个真正的低谷期。  
不过，这场“寒冬”也带来了清醒。它让研究者认识到，AI的发展必须更务实、更关注基础。知识表示、学习算法、计算复杂性等问题，开始被重新审视——这些反思，正是后来复苏的根基。

如果从“三种思维”的角度来总结这一段经历，可以说：数学思维与算法思维在理想蓝图上走得太快，而工程思维和现实复杂度的约束则给出了“冷静剂”。这种张力，最终推动了思维从“全靠规则”向“需要学习”的方向转向。

\section{专家系统与知识工程：AI的第二次兴起}

1970年代末至1980年代，人工智能迎来了第二次复兴。  
这一时期的核心成就是“专家系统”的兴起——它们通过模拟人类专家的知识与推理过程，在特定领域中展现出令人瞩目的问题解决能力。  
这也意味着AI第一次从实验室的理论研究，真正走向了可落地的实际应用。

在这一阶段，工程思维明显上升为主角：研究者开始更关注“在哪些领域真正好用”“如何把系统部署在真实环境”，同时也催生出一门新的技术艺术——知识工程。算法和数学不再单独领衔，而是为工程目标服务。

\subsection{专家系统的诞生与成功}

所谓“专家系统”，是一类以知识为核心的人工智能系统。它通常由两个部分构成：知识库与推理机。  
知识库存储领域专家积累的经验与规律，推理机则利用这些知识进行推断与决策。  
与早期依赖通用算法的AI不同，专家系统提出了一个重要理念——“知识比算法更重要”。这标志着AI研究的重心从“如何算”转向“知道什么”。

MYCIN 是最早也是最著名的专家系统之一。  
它由斯坦福大学于 1970 年代开发，用于辅助医生诊断细菌感染并推荐抗生素治疗方案。  
系统的知识库包含约 450 条规则，每条规则以“If … then …”形式表达。推理机会根据患者症状与检验结果匹配这些规则，从而推断可能的病原体并给出治疗建议。  
在当时的临床测试中，MYCIN 的诊断水平已接近人类专家，展示了 AI 在医学领域的巨大潜力。  
虽然 MYCIN 最终未能大规模投入临床使用，但它的设计理念和系统结构，成为后来专家系统的奠基蓝图。

另一个重要的成功案例是 DENDRAL。  
由爱德华·费根鲍姆（Edward Feigenbaum）等人领导开发，DENDRAL 专注于化学结构分析，能够根据实验数据自动推断分子的可能结构。  
这一系统显著提升了科研效率，也验证了一个关键观点：当AI系统获得专家级知识时，它可以在特定领域内展现出与人类专家相当的能力。

专家系统的成功，也催生了一个新的研究方向——知识工程。  
它关注的是：如何高效地获取、表示并应用知识。  
然而，在这一过程中，研究者发现最困难的部分并非推理，而是“如何获得专家知识”。  
这种困难被称为“知识获取瓶颈”。专家系统的性能，很大程度上取决于知识库的质量与推理机制的精度——也就是说，系统的智能，取决于它“知道得多准”。

\subsection{商业化与第二次AI寒冬}

专家系统的成功，让AI第一次真正走向商业世界。  
企业纷纷投入开发各类专家系统，广泛应用于医学诊断、金融分析、设备维护等领域。  
AI 不再只是实验室的学术实验，而开始创造可见的经济价值。  
到了 1980 年代，AI 被视为产业升级的关键引擎，一场前所未有的投资热潮由此展开。

但随着规模扩大，问题也随之而来。  
专家系统的开发与维护成本居高不下，知识获取缓慢，且难以应对环境的变化。  
最关键的是，它们缺乏自我学习能力——一旦知识库过时或出现错误，性能便会骤降。  
许多企业在投入巨资后未能获得预期回报，对AI的热情迅速冷却，信心再度受挫。

1987年，AI 行业再次步入寒冬。  
专用 AI 硬件（如 Lisp 机器）的市场崩溃成为转折点，投资者纷纷撤资，大量公司倒闭。  
这一波退潮让研究者意识到：AI 的商业化之路并非坦途。技术的边界与经济的现实，始终共同决定着智能发展的节奏。

从三种思维的角度来看，这一时期是工程思维“强势出场，又被现实教育”的过程：我们学会了如何让AI在特定场景下好用，也意识到仅靠工程堆叠、缺乏自我学习的系统缺乏生命力。这为之后“让机器自己学”的方向积蓄了动力。

\section{连接主义的复兴与深度学习的前夜}

在经历两次“寒冬”之后，AI研究者重新回到了一个根本问题——智能究竟从何而来？  
不同于符号主义对逻辑与规则的依赖，连接主义（Connectionism）提出了全新的视角：智能并非源于少数复杂规则，而是由大量简单单元的相互作用涌现出来的。  
神经网络模型正是在这种思想下重新焕发生机。

在这里，算法思维发生了一次重要转折：从“写规则+做推理”变成“设结构+靠学习”。数学思维通过逼近定理等结果为这种结构提供理论底气，而工程思维则还在等待算力与数据的成熟。

\subsection{神经网络的重生}

20世纪80年代，计算能力的提升与反向传播算法（backpropagation）的出现，为神经网络注入了新的生命。  
该算法通过计算输出误差对各层权重的梯度，逐层调整参数，使多层网络的训练成为可能。  
这不仅是算法上的突破，更意味着AI终于能在实践中处理更复杂的模式识别与非线性问题——神经网络开始显露真正的潜力。

反向传播算法的提出，使多层前馈神经网络不再停留在理论。  
1986年，大卫·鲁梅尔哈特（David Rumelhart）、杰弗里·辛顿（Geoffrey Hinton）与罗纳德·威廉姆斯（Ronald Williams）发表论文，系统阐述了该算法的原理与应用。  
这篇论文成为连接主义复兴的里程碑，让神经网络重新回到学术主舞台，并引发了新一轮研究热潮。

当然，这场复兴并非没有阻力。  
神经网络在当时仍面临多重挑战：计算资源匮乏、训练样本有限、模型可解释性不足。  
多层网络的训练计算量极大，而1980年代的硬件难以承受；数据稀缺，也使模型难以泛化。  
此外，部分研究者质疑神经网络无法像符号系统那样清晰表达知识。  
尽管如此，连接主义的探索仍然持续推进——这些努力，正是后来深度学习崛起前最宝贵的积累。

\subsection{从连接到学习的思想转折}

连接主义的核心，不在于如何编程，而在于如何让系统“自己学会”。  
符号主义依靠专家手工制定规则；连接主义则依赖算法，从数据中自动归纳出规律。  
这一思想转变意义深远——它标志着AI研究的焦点，从“获取知识”逐步转向“学习模型”。

这种转变不仅改变了AI的研究方法，也改变了人们对智能的看法。  
智能不再是一个被提前设定好的规则集合，而是一种能在学习中不断自我调整的动态过程。  
机器不再只是被动执行命令，而是主动在数据中发现模式、积累经验。  
“从编程到学习”——这一思想，成为现代AI最具标志性的特征之一。

连接主义的复兴，为后来深度学习的崛起铺平了道路。  
深度学习中的核心思想——分层特征表示、误差反向传播、多层神经结构——几乎都能追溯到这一时期的研究。  
可以说，没有连接主义的再次觉醒，就不会有后来的深度学习革命。

\section{深度学习的崛起：从经验到智能}

进入21世纪后，AI迎来了第三次真正意义上的复兴。  
随着计算能力的爆发式提升、大规模数据的积累以及算法的持续改进，一场以“深度学习（Deep Learning）”为核心的智能革命开始了。  
深度学习以多层神经网络为基础，能够自动从数据中学习特征与模式，并在语音识别、图像识别、自然语言处理等领域取得了前所未有的突破。

深度学习的核心思想在于“多层抽象”。  
它通过层层非线性变换，使机器能够从原始数据中自动提取有用信息。  
在深度学习出现之前，传统机器学习严重依赖人工特征提取——例如图像识别中需要人工设计边缘检测、角点检测等特征。  
如今，这些工作被交给神经网络自动完成，大大减少了人工干预，也使模型能在更复杂的数据中发现隐藏结构。

2006年，杰弗里·辛顿（Geoffrey Hinton）等人提出了深度信念网络（Deep Belief Network, DBN），  
这被普遍认为是深度学习时代的起点。  
DBN 通过“逐层无监督预训练 + 有监督微调”的方式，有效解决了深层网络难以收敛的问题。  
随后，卷积神经网络（CNN）与循环神经网络（RNN）的兴起，使深度学习在结构表达与性能上进一步跃升。

2012年，AI历史迎来了决定性的一刻。  
辛顿的学生亚历克斯·克里热夫斯基（Alex Krizhevsky）在 ImageNet 图像识别竞赛中，  
凭借深度卷积神经网络（AlexNet）取得了远超传统算法的成绩——错误率降低超过一半。  
这一结果震动了学术界与产业界，标志着深度学习的全面崛起。  
自此之后，深度学习成为AI研究的主流方向，也重新定义了“机器能看、能听、能理解”的边界。

深度学习的成功并非偶然。  
它依托于三个关键支点：算法的改进、计算能力的跃升与大数据的积累。  
GPU 的引入让神经网络的训练速度提升数十倍；大规模标注数据集（如 ImageNet）为模型提供了丰富的学习素材；  
同时，算法层面的创新——如 ReLU 激活函数与 Dropout 技术——显著提升了模型的稳定性与泛化能力。  
三者合力，共同催生了AI的第三次浪潮。

深度学习的影响远不止于技术突破。  
它在哲学层面上，也改变了人们对智能的理解。  
深度学习展示了另一种智能的路径：  
不是依赖人类事先制定的规则，而是通过经验的积累与自我调整获得能力。  
机器不再依赖固定的知识体系，而是在数据中不断修正自身结构——  
这种“从经验中学习”的能力，让AI的学习方式首次真正接近人类。

从这个意义上说，深度学习不仅是一种算法，更是一种新的思维范式。  
它提醒我们：智能的核心，不在于拥有完备的规则，而在于在不确定环境中自我调整的能力。  
这种思维的转变，使人工智能进入了一个以“学习”为核心的新阶段。

如果把这一阶段放回“三种思维”的框架中来看，会发现：数学思维通过概率、优化和信息论为模型提供理论支撑；算法思维通过BP、CNN、RNN乃至Transformer构造出一整套可训练的结构；工程思维则借助GPU、分布式系统与大规模数据，把这些模型带到了现实世界的规模上。三种思维第一次形成了真正意义上的“大合唱”。

\section{小结：AI的思维演化之路}

回顾人工智能的发展，我们会发现我们会发现，AI的演化并非直线前进，而更像是一个不断循环、自我修正的螺旋过程。 
从符号主义的逻辑推理，到连接主义的神经计算，再到深度学习的经验学习，AI经历了多次关键的思想转向。  而这些转向，从未是简单的彼此取代，而是思维范式的层层叠加，共同构筑出今日AI的广阔图景。

如果用本章开头提出的三种思维来概括：  
符号主义强调逻辑与规则，揭示了智能的形式结构；  
连接主义关注神经与关联，探索了智能的生理机制；  
深度学习聚焦经验与自适应，展示了智能的学习本质。  

符号主义阶段，数学思维与算法思维主导舞台，试图用逻辑与程序穷尽智能的结构；  
专家系统与知识工程阶段，工程思维走到前台，强调“在具体领域真正好用”；  
连接主义与深度学习阶段，则是在数学与算法的支撑下，让“从经验中学习”的能力真正落地。  
三种思维各有局限，却也在相互借鉴中不断成长，正是在这种互补中，AI逐渐走向成熟。

从更宏观的视角来看，AI 的发展正体现了科学思维的普遍规律：  
理论推动实践，实践反哺理论；危机之中，往往孕育着新的机遇。  
每一次“瓶颈期”，都成为思维更新的催化剂。  
正是在这些困惑与反思的时刻，人工智能一步步跨越自身的边界。

因此，理解AI的历史，不只是追溯技术的演进，更是在体会一种思维的成长。  
数学思维让我们能够“定义智能”，算法思维让我们能够“求解智能”，工程思维让我们能够“实现智能”。  
当这三种思维相互作用，我们才真正触及了人工智能的核心所在——一种在人类探索中持续进化的智慧形态。

AI的发展历程启示我们：智能并非被“制造”出来的产物，而是在不断探索中被“发现”的过程。  
每一次理论的突破、每一次技术的跃迁，都是人类理解自身思维方式的一次深化。  
而人工智能的未来，也必将在这条思维演化的道路上，继续前行。

